\begin{notechapter}
\chapter{Estimators and Sampling Distributions}\label{ch:stats-and-proba}
\chaptersubtitle{What is fixed, what is random, and how good is an estimate?}

\begin{chapterintro}
Every number a learning algorithm reports --- a fitted weight, a training error,
a test accuracy --- is computed from random data, and would have come out
differently with another dataset. This chapter gives the vocabulary and the
basic tools to reason about that variability. You will learn
\begin{itemize}[nosep]
  \item to distinguish an \emph{unknown} quantity from a \emph{random} one, and an
        \emph{estimator} from an \emph{estimate};
  \item what the \emph{sampling distribution} of an estimator is;
  \item the bias, variance and mean squared error of an estimator, and the
        decomposition linking them;
  \item why the sample variance divides by $n-1$, and why an unbiased estimator
        is not always the best one;
  \item how quantities that depend on the whole training set are studied, and
        how the bootstrap approximates their distribution from a single dataset.
\end{itemize}
\prerequisites{expectation and variance; \cref{learning-problem:sec:sample}
(the training sample, and the meanings of ``sample'').}
\end{chapterintro}

\section{Unknown quantities and random quantities}\label{stats-and-proba:sec:unknown-random}

Statistics deals with two kinds of quantities that are easily confused,
because we are uncertain about both.

\begin{description}
  \item[An unknown quantity] is a fixed number that we do not know: the mean
        height of a population, the true slope $\beta_1$ of a regression model,
        the risk $\Risk(f)$ of a predictor. It does not vary; our goal is to
        \emph{estimate} it.
  \item[A random quantity] varies from one repetition of an experiment to the
        next according to a probability distribution: the outcome of a coin
        toss, the noise $\varepsilon$ in a regression, the empirical risk
        computed on a random sample. Our goal is to \emph{describe its
        distribution}.
\end{description}

\begin{table}[ht]
\centering
\caption{Unknown versus random quantities.}
\label{stats-and-proba:tab:unknown-random}
\begin{tabular}{@{}lll@{}}
\toprule
 & unknown quantity & random quantity \\
\midrule
nature & fixed, not known & varies from draw to draw \\
described by & a single (unknown) value & a probability distribution \\
goal & estimate it & model its variability \\
example & true parameter $\beta_1$, risk $\Risk(f)$ & coin toss, noise $\varepsilon$, $\ERisk(f)$ \\
\bottomrule
\end{tabular}
\end{table}

\begin{pitfall}[Uncertain is not the same as random]
In the frequentist view adopted in this chapter, a parameter is unknown but
\emph{not} random: it makes no sense to speak of ``the probability that
$\beta_1 > 0$''. Probabilities describe the data and the quantities computed
from them. The Bayesian view (\cref{ch:bayesian-reasoning}) deliberately treats
unknown parameters as random variables, in order to express uncertainty about
them with probabilities. Both views are coherent; mixing them is what causes
confusion.
\end{pitfall}

\section{Estimators and their sampling distribution}\label{stats-and-proba:sec:estimators}

Let $X_1,\dots,X_n$ be an i.i.d.\ sample from a distribution that depends on
an unknown parameter $\theta$ (for instance its mean). As everywhere in this book, ``sample'' means the whole
collection (\cref{learning-problem:sec:sample-meanings}); each $X_i$ is an
observation.

\begin{definition}[Estimator, estimate, sampling distribution]
An \emph{estimator} of $\theta$ is a function of the sample,
$\hat\theta = T(X_1,\dots,X_n)$; it is a random variable. Its value
$T(x_1,\dots,x_n)$ on an observed dataset is an \emph{estimate}, a number. The
probability distribution of $\hat\theta$ --- the distribution of the values it
would take over all possible samples --- is its \emph{sampling distribution}.
\end{definition}

The sample mean $\bar X_n = \frac1n\sum_iX_i$ is an estimator of the mean $\mu$;
$\bar x = 5.2$ computed on a dataset is an estimate. The empirical risk of
\cref{learning-problem:def:emprisk} is an estimator of the risk.

\begin{example}[An exact sampling distribution]\label{stats-and-proba:ex:binomial}
A classifier with error rate $p$ is tested on $n$ independent examples. The
number of errors is binomial, so the estimated error rate $\hat p = K/n$ takes
the values $0, \frac1n, \dots, 1$ with $\Prob(\hat p = k/n) = \binom nk p^k(1-p)^{n-k}$.
With $p = 0.1$ and $n = 50$, $\Prob(\hat p\leq 0.04) = \Prob(K\leq2)\approx0.11$:
about one test set in nine would make this classifier look more than twice as
good as it is.
\end{example}

In general the sampling distribution is unknown, because it depends on the
unknown distribution of the data. It can be approximated by limit theorems
(the central limit theorem, \cref{ch:lln-clt-inference}) or by resampling (the
bootstrap, \cref{stats-and-proba:sec:bootstrap}).

\section{Bias, variance and mean squared error}\label{stats-and-proba:sec:bias-variance}

\begin{definition}[Bias, variance, MSE]
The \emph{bias}, \emph{variance} and \emph{mean squared error} of an estimator
$\hat\theta$ of $\theta$ are
\[
  \mathrm{Bias}(\hat\theta) = \E[\hat\theta] - \theta,\qquad
  \Var(\hat\theta) = \E\bigl[(\hat\theta - \E\hat\theta)^2\bigr],\qquad
  \mathrm{MSE}(\hat\theta) = \E\bigl[(\hat\theta - \theta)^2\bigr].
\]
$\hat\theta$ is \emph{unbiased} if its bias is zero.
\end{definition}

The bias depends on the unknown $\theta$ and the variance on the unknown
distribution of the data, so neither is known in practice; but both can be
studied mathematically, and the variance can be estimated from the data.

\begin{proposition}[MSE decomposition]\label{stats-and-proba:prop:mse}
$\mathrm{MSE}(\hat\theta) = \mathrm{Bias}(\hat\theta)^2 + \Var(\hat\theta)$.
\end{proposition}

\begin{proof}
Write $\hat\theta - \theta = (\hat\theta - \E\hat\theta) + (\E\hat\theta - \theta)$.
The second term is a constant, and the first has mean zero, so the cross term
vanishes when the square is expanded and its expectation taken.
\end{proof}

This is the estimator-level version of the bias--variance decomposition of
\cref{erm:prop:bias-variance}: an estimator can be wrong on average (bias) or
unstable from sample to sample (variance), and its MSE adds the two.

\begin{intuition}[Darts]
Think of repeated samples as repeated throws at a target located at $\theta$.
The bias is the distance from the centre of the cloud of hits to the target;
the variance is the spread of the cloud. A tight cloud off-centre and a wide
cloud centred on the target can have the same mean squared distance.
\end{intuition}

\subsection{The sample mean}

\begin{proposition}\label{stats-and-proba:prop:mean}
If the $X_i$ have mean $\mu$ and variance $\sigma^2$, then
$\E[\bar X_n] = \mu$ and $\Var(\bar X_n) = \sigma^2/n$.
\end{proposition}

\begin{proof}
Linearity of expectation gives the mean; independence gives
$\Var(\sum_iX_i) = \sum_i\Var(X_i) = n\sigma^2$, and dividing by $n^2$ gives the
variance.
\end{proof}

The \emph{standard error} $\sigma/\sqrt n$ of the mean decreases like $1/\sqrt n$:
four times more data halves the uncertainty. It is estimated by $s/\sqrt n$,
with $s$ the sample standard deviation defined next.

\subsection{The sample variance and Bessel's correction}

\begin{proposition}[Bessel's correction]\label{stats-and-proba:prop:bessel}
Let $S^2 = \frac1{n-1}\sum_{i=1}^n(X_i - \bar X_n)^2$. Then $\E[S^2] = \sigma^2$,
whereas $\hat\sigma^2 = \frac1n\sum_i(X_i - \bar X_n)^2 = \frac{n-1}{n}S^2$ has
$\E[\hat\sigma^2] = \frac{n-1}n\sigma^2$.
\end{proposition}

\begin{proof}
Centring at the true mean, $X_i - \bar X_n = (X_i - \mu) - (\bar X_n - \mu)$, and
\[
  \sum_i(X_i - \bar X_n)^2 = \sum_i(X_i - \mu)^2 - n(\bar X_n - \mu)^2,
\]
because the cross term is $-2(\bar X_n - \mu)\sum_i(X_i - \mu) = -2n(\bar X_n - \mu)^2$.
Taking expectations, the first sum has mean $n\sigma^2$ and the second term
$n\Var(\bar X_n) = \sigma^2$, so $\E\sum_i(X_i - \bar X_n)^2 = (n-1)\sigma^2$.
\end{proof}

\begin{intuition}[Why $n - 1$?]
The deviations are measured from $\bar X_n$, which is computed from the same
data and is therefore closer to them than the true mean $\mu$ is: the sum of
squared deviations is smallest around $\bar X_n$. Dividing by $n$ would
underestimate the spread. Another reading: the $n$ deviations
$X_i - \bar X_n$ sum to zero, so only $n-1$ of them are free --- one
``degree of freedom'' has been spent estimating the mean.
\end{intuition}

\begin{example}[Five measurements]\label{stats-and-proba:ex:five}
Heights $170, 165, 180, 175, 160$ (cm) have $\bar x = 170$ and squared
deviations $0, 25, 100, 25, 100$, summing to $250$. The unbiased estimate is
$s^2 = 250/4 = 62.5$ ($s\approx7.9$ cm) and the maximum-likelihood estimate
under a Gaussian model is $\hat\sigma^2 = 250/5 = 50$ ($\hat\sigma\approx7.1$ cm).
The standard error of the mean is estimated by $s/\sqrt5\approx3.5$ cm.
\end{example}

\begin{pitfall}[Unbiased is not the same as best]
\begin{itemize}[nosep]
  \item $S$ is a \emph{biased} estimator of $\sigma$, even though $S^2$ is
        unbiased for $\sigma^2$: since the square root is concave, Jensen's
        inequality gives $\E[S]\leq\sqrt{\E[S^2]} = \sigma$. Unbiasedness does
        not survive nonlinear transformations.
  \item For Gaussian data, $\Var(S^2) = 2\sigma^4/(n-1)$, so
        $\mathrm{MSE}(S^2) = 2\sigma^4/(n-1)$, while the biased $\hat\sigma^2$
        has $\mathrm{MSE}(\hat\sigma^2) = \sigma^4/n^2 + 2\sigma^4(n-1)/n^2 = (2n-1)\sigma^4/n^2$,
        which is \emph{smaller}. Accepting a little bias to reduce variance is
        exactly what regularization does (\cref{erm:sec:bias-variance}).
\end{itemize}
\end{pitfall}

\section{Quantities that depend on the training set}\label{stats-and-proba:sec:training-set}

In machine learning the estimator of interest is often the output of a whole
learning procedure. The training sample
$D_n = ((X_1,Y_1),\dots,(X_n,Y_n))$ is random (\cref{learning-problem:def:sample}),
so the trained model $\hat f_{D_n}$, its parameters, and its test error are
random as well. A typical object of study is an expectation over training sets,
\[
  \E_{D_n}\bigl[\phi(D_n)\bigr] = \int\phi(d)\,P^n(\mathrm{d}d),
\]
for instance $\phi(D_n) = \Risk(\hat f_{D_n})$, the expected risk of the algorithm
(\cref{erm:sec:model-vs-algorithm}). Its integrand involves the unknown
distribution $P^n$, so it cannot be computed directly.

\paragraph{Monte Carlo over datasets.} If we could draw $m$ independent training
sets $D^{(1)},\dots,D^{(m)}$ from the population, the average
$\frac1m\sum_k\phi(D^{(k)})$ would estimate $\E_{D_n}[\phi(D_n)]$ without bias,
with an error of order $1/\sqrt m$ (\cref{ch:monte-carlo}). This is how such
quantities are studied in simulations, where the data-generating distribution is
known.

\subsection{The bootstrap}\label{stats-and-proba:sec:bootstrap}

With real data we have a single dataset. The bootstrap (Efron, 1979) replaces
the unknown distribution $P$ by the \emph{empirical distribution} $\hat P_n$ of
the data, which puts mass $\frac1n$ on each observed point, and samples from it.

\begin{center}
\begin{minipage}{0.9\linewidth}
\textbf{Bootstrap estimate of the standard error of a statistic $T$.}
\begin{enumerate}[nosep]
  \item For $b = 1,\dots,B$: draw $n$ points \emph{with replacement} from the
        dataset, giving a resample $d^{*(b)}$, and compute $T^{*(b)} = T(d^{*(b)})$.
  \item Report the standard deviation of $T^{*(1)},\dots,T^{*(B)}$ as the
        estimated standard error of $T$. Their empirical distribution
        approximates the sampling distribution of $T$.
\end{enumerate}
\end{minipage}
\end{center}

\begin{keyidea}[What the bootstrap estimates]
The bootstrap computes (by Monte Carlo) the sampling distribution of $T$ when
the data come from $\hat P_n$ instead of $P$. It works when $\hat P_n$ is a good
stand-in for $P$ --- large $n$, and a statistic that depends smoothly on the
distribution --- and needs no formula for the standard error. A resample
contains on average only $1 - (1-\frac1n)^n\approx63\%$ of the distinct
original points.
\end{keyidea}

\begin{pitfall}[Bootstrap and error estimates]
Evaluating a model trained on a bootstrap resample on the \emph{original} data
is optimistic: about $63\%$ of the evaluation points were seen in training. Use
the points left out of each resample (the ``out-of-bag'' points) for
evaluation, or use cross-validation (\cref{measuring-performance:sec:cv}).
\end{pitfall}

\begin{chaptersummary}
\begin{itemize}[nosep]
  \item A parameter is fixed but unknown; an estimator is a random variable
        computed from the sample; an estimate is its value on the observed data.
  \item The sampling distribution of an estimator describes its values across
        samples. Its mean defines the bias, its spread the variance, and
        $\mathrm{MSE} = \mathrm{Bias}^2 + \Var$.
  \item $\bar X_n$ is unbiased with variance $\sigma^2/n$; $S^2$ with divisor
        $n-1$ is unbiased for $\sigma^2$, but unbiasedness is not preserved by
        nonlinear maps and is not the same as small MSE.
  \item Trained models and their errors are functions of the random training
        set; expectations over training sets are approximated by Monte Carlo
        when data can be simulated, and by the bootstrap from a single dataset.
\end{itemize}
\end{chaptersummary}

\section*{Exercises}
\addcontentsline{toc}{section}{Exercises}

\begin{exercise}\label{stats-and-proba:exr:1}
Compute $\Prob(\hat p\leq0.04)$ in \cref{stats-and-proba:ex:binomial} exactly,
and compare with the normal approximation
$\hat p\approx\mathcal N(p, p(1-p)/n)$.
\solutionlink{stats-and-proba:exr:1}
\end{exercise}

\begin{exercise}\label{stats-and-proba:exr:2}
Let $X_1,\dots,X_n$ be uniform on $[0,\theta]$. Show that $2\bar X_n$ is an
unbiased estimator of $\theta$ and compute its MSE. Then show that
$M = \max_iX_i$ satisfies $\E[M] = \frac n{n+1}\theta$, and compare the MSE of
$\frac{n+1}nM$ with that of $2\bar X_n$.
\hint{$\Prob(M\leq x) = (x/\theta)^n$ for $0\leq x\leq\theta$.}
\solutionlink{stats-and-proba:exr:2}
\end{exercise}

\begin{exercise}\label{stats-and-proba:exr:3}
Show that the probability that a given point appears in a bootstrap resample of
size $n$ is $1 - (1 - \frac1n)^n$, and that it tends to $1 - e^{-1}\approx0.632$.
\solutionlink{stats-and-proba:exr:3}
\end{exercise}

\begin{exercise}\label{stats-and-proba:exr:4}
Among estimators of the form $cS^2$ of $\sigma^2$ for Gaussian data, find the
constant $c$ that minimizes the MSE.
\hint{use $\Var(S^2) = 2\sigma^4/(n-1)$.}
\solutionlink{stats-and-proba:exr:4}
\end{exercise}
\end{notechapter}
